%% ============================================================================
%% プリアンブル。main.tex の \documentclass{acmart} の「後」に \input される。
%%
%% 以下は acmart.cls v2.20 が既に読み込み済み。重複して \usepackage しないこと。
%% またここでオプションを渡すこともできない（既に読まれた後なので手遅れ）。
%% オプションが要る場合は main.tex の \documentclass より前で
%% \PassOptionsToPackage を使う。
%%
%%   amssymb  balance  booktabs  caption  comment  fancyhdr  float  fontenc
%%   geometry  graphicx  hyperref  libertine  microtype  natbib  newtxmath
%%   pifont  setspace  xcolor
%%
%% したがって \checkmark と \mathbb は amssymb、\toprule は booktabs、
%% \ding{51} は pifont、\textcolor は xcolor から、そのまま使える。
%% ============================================================================


%% --- 数式 -------------------------------------------------------------------
\usepackage{mathtools}   % 数式環境一式（amsmath を含む）、\coloneqq など
\usepackage{bm}          % 太字の数式 \bm{x}


%% --- 図 ---------------------------------------------------------------------
\usepackage{subcaption}  % subfigure / subtable 環境
\usepackage{placeins}    % \FloatBarrier で図の流れ込みを止める

\graphicspath{{figs/}{figures/}}

%% サブ図のラベルを (a) (b) 形式にする
\renewcommand{\thesubfigure}{(\alph{subfigure})}
\captionsetup[subfigure]{labelformat=simple}


%% --- 表 ---------------------------------------------------------------------
\usepackage{array}           % \newcolumntype、m{} / b{} 列
\usepackage{tabularx}        % 指定幅まで伸びる X 列
\usepackage{multirow}        % \multirow でセルの縦結合
\usepackage{makecell}        % セル内改行 \makecell{a\\b}
\usepackage{colortbl}        % \rowcolor / \cellcolor（xcolor は読み込み済み）
\usepackage{threeparttable}  % 表幅に揃う脚注
\usepackage{adjustbox}       % \adjustbox{max width=\linewidth}{...}、回転

\definecolor{lightgray}{gray}{0.85}

%% 幅を固定した左寄せ・中央・右寄せの列。tabularx と組み合わせて使う
\newcolumntype{L}[1]{>{\raggedright\arraybackslash}p{#1}}
\newcolumntype{C}[1]{>{\centering\arraybackslash}p{#1}}
\newcolumntype{R}[1]{>{\raggedleft\arraybackslash}p{#1}}


%% --- 数値と単位 -------------------------------------------------------------
\usepackage{siunitx}     % \num{1234567} -> 1,234,567、\SI{3.5}{\milli\second}
\sisetup{
  detect-all,            % 周囲のフォントに追従させる
  group-separator = {,},
  table-align-uncertainty = true,
}


%% --- 箇条書き ---------------------------------------------------------------
%% itemize / enumerate / description の余白や記号を調整する
\usepackage{enumitem}
\setlist{leftmargin=*, itemsep=2pt, topsep=2pt, parsep=0pt}


%% --- アルゴリズム -----------------------------------------------------------
\usepackage{algorithm}      % algorithm 環境（図表と同じくフロートになる）
\usepackage{algpseudocode}  % \State \If \While などの擬似コード


%% --- URL とリンク -----------------------------------------------------------
\usepackage{xurl}        % 長い URL を任意の位置で折り返す

\hypersetup{
  colorlinks = true,
  linkcolor  = black,
  urlcolor   = black,
  citecolor  = black,
  breaklinks = true,
}


%% --- テキスト ---------------------------------------------------------------
\usepackage{xspace}      % \newcommand{\sys}{FooBar\xspace} で後続の空白を保つ


%% --- 相互参照 ---------------------------------------------------------------
%% \cref / \Cref で参照先の種類（Figure, Table ...）を自動で補う。
%% 上のパッケージの参照コマンドを書き換えるので、必ず最後に読み込む。
%%   capitalise -> 「figure 1」ではなく「Figure 1」
%%   noabbrev   -> 「Fig. 1」「Sec. 2」ではなく「Figure 1」「Section 2」
%%   nameinlink -> 「Figure」の語もリンクに含める（数字だけをリンクにしない）
\usepackage[capitalise, noabbrev, nameinlink]{cleveref}

\crefname{subfigure}{Figure}{Figures}
\Crefname{subfigure}{Figure}{Figures}
\crefname{equation}{Equation}{Equations}
\Crefname{equation}{Equation}{Equations}
\crefname{appendix}{Appendix}{Appendices}
\Crefname{appendix}{Appendix}{Appendices}
\crefname{algorithm}{Algorithm}{Algorithms}
\Crefname{algorithm}{Algorithm}{Algorithms}


%% --- 著者定義のマクロ -------------------------------------------------------
\newcommand{\yes}{\ensuremath{\checkmark}}
\newcommand{\partly}{\ensuremath{\triangle}}
\newcommand{\no}{\textemdash}

%% 執筆中のメモ。\drafttrue を \draftfalse にすれば一括で消える
\newif\ifdraft
\drafttrue

\ifdraft
  \newcommand{\todo}[1]{\textcolor{red}{[TODO: #1]}}
  \newcommand{\note}[2]{\textcolor{blue}{[#1: #2]}}
\else
  \newcommand{\todo}[1]{}
  \newcommand{\note}[2]{}
\fi
